\documentclass{article}
\usepackage{booktabs,longtable,geometry,array,color}
\geometry{margin=1in}
\begin{document}

\section*{M019: Radioactive Iodine (RAI) Outcomes Analysis}
\textit{Generated 2026-05-05 | N=862 RAI recipients | Malignant cohort N=4,019}

\bigskip
\noindent\textbf{Critical Data Limitations:}
\begin{itemize}
  \item 75.2\% of RAI recipients lack dose data (648/862); only 35 patients have confirmed-with-dose tier
  \item 88.5\% have NULL or unknown RAI intent (763/862: 613 NULL + 150 unknown)
  \item 95.9\% lack formal ATA response category (827/862)
  \item TgAb interference affects 58 (6.7\%) patients --- Tg unreliable in these patients
  \item 34.3\% lack Tg trajectory data (296/862)
\end{itemize}

% -------------------------------------------------------------------
\subsection*{Table 1: RAI Utilization by AJCC Stage (Malignant Cohort, N=4,019)}
\begin{tabular}{lrrr}
\toprule
AJCC Stage & Total N & Received RAI & \% RAI \\
\midrule
I   & 1,537 & 192 & 12.5\% \\
II  & 1,651 & 311 & 18.8\% \\
III &     9 &   1 & 11.1\% \\
IVB &   816 & 146 & 17.9\% \\
\midrule
\textit{Overall} & 4,019 & 650 & 16.2\% \\
\bottomrule
\end{tabular}

\vspace{0.3cm}
\noindent\textit{Note: RAI receipt in the full cohort (malignant only, CPM) = 650/4,019 (16.2\%).
The analytic view (m019\_rai\_outcomes\_analytic\_v1) contains 862 patients using a broader reconciliation definition.
No Stage IVA patients in current cohort; AJCC 8th Ed used throughout.}

% -------------------------------------------------------------------
\subsection*{Table 2: RAI Receipt by Histology (Top Categories)}
\begin{tabular}{lrrr}
\toprule
Histology & Total N & Received RAI & \% RAI \\
\midrule
PTC         & \multicolumn{3}{c}{\textit{see rai\_utilization\_patterns.csv}} \\
FTC         & & & \\
HCC         & & & \\
PDTC        & & & \\
\bottomrule
\end{tabular}

\vspace{0.3cm}
\noindent\textit{Full breakdown available in rai\_utilization\_patterns.csv}

% -------------------------------------------------------------------
\subsection*{Table 3: RAI Temporal Trends (De-escalation)}
\begin{tabular}{lrrr}
\toprule
Year & Malignant N & Received RAI & \% RAI \\
\midrule
2010 &  85 & 11 & 12.9\% \\
2013 & 169 & 23 & 13.6\% \\
2015 & 171 & 40 & 23.4\% \\
2016 & 162 & 49 & 30.2\% \\
2017 & 225 & 58 & 25.8\% \\
2018 & 220 & 53 & 24.1\% \\
2019 & 328 & 118 & \textbf{36.0\%} \\
2020 & 251 & 33 & 13.1\% \\
2021 & 286 & 31 & 10.8\% \\
2022 & 279 & 43 & 15.4\% \\
2023 & 369 & 44 & 11.9\% \\
2024 & 379 & 39 & 10.3\% \\
\bottomrule
\end{tabular}

\vspace{0.3cm}
\noindent\textit{2019 represents the peak RAI utilization year (36.0\%). Post-2019 de-escalation
to $\sim$10--15\% consistent with ATA 2015 guideline uptake and growing low-risk management evidence.}

% -------------------------------------------------------------------
\subsection*{Table 4: Recurrence by Tg Trajectory Class (N=509, excluding TgAb interference)}
\textit{Suppressed = reference category. Excludes 57 TgAb-interference patients.}

\begin{tabular}{lrrrrr}
\toprule
Tg Trajectory & N & Recurred & \% & OR (95\% CI) & p-value \\
\midrule
Suppressed (ref)   & 251 & 57 & 22.7\% & 1.00 (ref) & --- \\
Low stable         & 102 & 25 & 24.5\% & 1.11 (0.64--1.89) & 0.82 \\
Detectable stable  &  82 & 18 & 22.0\% & 0.96 (0.53--1.75) & 1.00 \\
Rising             &  71 & 23 & 32.4\% & 1.63 (0.91--2.91) & 0.13 \\
\bottomrule
\end{tabular}

\vspace{0.3cm}
\noindent\textit{Post-RAI Tg nadir ROC AUC = 0.509 (optimal threshold 0.60 ng/mL: Sens=0.31, Spec=0.75).
The weak discriminative value reflects the overall high recurrence in this surgically-selected cohort,
limited Tg data coverage (59\% of RAI cohort), and residual confounding by disease severity.}

% -------------------------------------------------------------------
\subsection*{Table 5: Recurrence by Dose Category (N=214 with dose data)}
\textit{CAUTION: Only 214/862 (24.8\%) have dose data. No patients met ``confirmed\_with\_dose'' confidence tier.
Interpret with extreme caution --- higher doses reflect higher-risk disease, not causal dose effects.}

\begin{tabular}{lrrrr}
\toprule
Dose Category & N & Median Dose (mCi) & Recurred & \% \\
\midrule
$<$30 mCi (low)              &  3 &  10.5 &  0 &  0.0\% \\
30--100 mCi (standard)       & 16 &  74.2 &  4 & 25.0\% \\
100--150 mCi (high)          & 98 & 124.4 & 26 & 26.5\% \\
$>$150 mCi (therapeutic)     & 97 & 176.0 & 32 & 33.0\% \\
\bottomrule
\end{tabular}

\vspace{0.3cm}
\noindent\textit{Logistic regression (log-dose vs recurrence): OR per log-unit = 2.93 (p=0.019).
This likely reflects indication bias --- patients receiving higher doses had more advanced disease.}

% -------------------------------------------------------------------
\subsection*{Table 6: RAI vs No-RAI Comparison (Malignant Cohort)}
\textit{RAI N=650 vs No-RAI N=3,369. Note: RAI recipients are a \textbf{selected, higher-risk} subgroup.}

\begin{tabular}{lrrr}
\toprule
Variable & RAI (N=650) & No-RAI (N=3,369) & p-value \\
\midrule
Age (mean $\pm$ SD)    & 47.8 $\pm$ 15.9 & 51.2 $\pm$ 15.6 & $<$0.001 \\
Tumor size cm (mean)   & 2.47 $\pm$ 1.86 & 2.19 $\pm$ 1.97 & 0.001 \\
Female sex             & 73.8\% & 72.9\% & 0.66 \\
BRAF positive          & 11.1\% & 6.1\% & $<$0.001 \\
LN positive            & 68.9\% & 58.8\% & $<$0.001 \\
AJCC Stage I/II        & 77.3\% & 79.8\% & --- \\
AJCC Stage IVB         & 22.5\% & 19.9\% & --- \\
Recurrence (crude)     & 23.5\% & 10.4\% & $<$0.001 \\
\midrule
\multicolumn{4}{l}{\textit{IPTW-ATT (adjusted for age, tumor size, Stage III+, LN status):}} \\
Recurrence (RAI)       & 23.7\% & --- & --- \\
Recurrence (No-RAI)    & 11.8\% & --- & --- \\
Risk Difference (ATT)  & \multicolumn{2}{l}{+11.9 pp (RAI $>$ No-RAI)} & --- \\
\bottomrule
\end{tabular}

\vspace{0.3cm}
\noindent\textit{Propensity model: n=3,876 complete cases. IPTW-ATT partially adjusts for indication bias
but residual confounding is expected. PS overlap was limited (mean PS: RAI=0.175 vs No-RAI=0.157),
indicating RAI recipients have modestly higher predicted probability. Formal causal inference requires
richer PS specification with all ATA risk factors. The higher crude recurrence in RAI recipients is
consistent with selection of higher-risk patients for RAI treatment.}

% -------------------------------------------------------------------
\subsection*{Table 7: Timing from Surgery to First RAI (N=336 with timing data)}
\begin{tabular}{lrrrr}
\toprule
Timing Group & N & Median Days & Recurred & \% \\
\midrule
Early ($<$90 days)     &  69 &   60 & 15 & 21.7\% \\
Mid (90--180 days)     &  42 &  125 & 12 & 28.6\% \\
Late ($>$180 days)     & 225 & 2,234 & 46 & 20.4\% \\
\bottomrule
\end{tabular}

\vspace{0.3cm}
\noindent\textit{Median overall timing = not shown (missing in 61\% of RAI cohort).
The ``late'' group median of 2,234 days suggests many ``late'' patients received subsequent RAI
treatments for recurrence (not initial ablation). Timing analysis should be restricted to
first-treatment ablation intent patients; intent is missing in 88.5\% of cohort.}

% -------------------------------------------------------------------
\subsection*{Table 8: ATA Dynamic Risk Assessment (N=35 --- Descriptive Only)}
\textit{N too small for inferential statistics. Cross-tabulation only.}

\begin{tabular}{lrrr}
\toprule
ATA Response Category & N & Recurred & \% \\
\midrule
Excellent              &  1 &  0 &  0.0\% \\
Indeterminate          &  4 &  1 & 25.0\% \\
Biochemical incomplete &  6 &  1 & 16.7\% \\
Structural incomplete  & 23 & 11 & 47.8\% \\
Insufficient data      &  1 &  1 & 100\% \\
\bottomrule
\end{tabular}

\vspace{0.3cm}
\noindent\textit{Response categories are directionally consistent with ATA 2015 framework expectations
but the N=35 sample (4.1\% of RAI cohort) precludes statistical inference.}

% -------------------------------------------------------------------
\subsection*{Table 9: TgAb Interference}
\begin{tabular}{lrrrr}
\toprule
Group & N & \% of RAI cohort & Recurred & \% \\
\midrule
TgAb positive  & 58 & 6.7\% & 13 & 22.4\% \\
TgAb negative  & 510 & 59.2\% & 123 & 24.1\% \\
Unknown/missing & 294 & 34.1\% & --- & --- \\
\bottomrule
\end{tabular}

\vspace{0.3cm}
\noindent\textit{TgAb positivity does not significantly alter recurrence rate vs TgAb-negative patients.
However, Tg values are unreliable as response markers in TgAb-positive patients; their Tg trajectory
classifications should be interpreted with caution or excluded from Tg-based response analyses.}

\end{document}
